\chapter{Replication, ISR, and High Watermarks}
This chapter moves from a single-machine log to replication. \textbf{First, read the boundaries of the course's distributed model carefully:} three independent brokers each have a TCP server, threads, and a disk directory; a follower must fetch records over TCP, and cannot replicate by sharing memory. However, they run in the same JVM and are managed by one trusted ClusterAuthority. The authority linearizes metadata, epoch, and HW state, making teaching failovers controllable. It is not KRaft, has no controller consensus or authority HA, and makes no guarantee about cluster recovery after loss of the entire JVM/authority.

Replication is meaningful only when append, reporting, and visibility boundaries are consistent with one another. LEO is the next offset in each individual log. ISR is the set of replicas currently allowed to participate in acknowledgments. HW is the committed boundary readable by ordinary consumers; consumers can read only offsets < HW. Replicas fetch over a real network; the authority provides consistent state and role decisions, but does not move data.

\begin{figure}[ht]
\centering
\begin{tikzpicture}[font=\small,node distance=20mm]
\node[draw,rounded corners,minimum width=31mm,minimum height=10mm] (leader) {leader log, LEO=5};
\node[draw,rounded corners,minimum width=31mm,minimum height=10mm,right=of leader] (follower) {follower log, LEO=3};
\node[draw,rounded corners,minimum width=33mm,minimum height=10mm,below=of $(leader)!0.5!(follower)$] (auth) {authority: ISR, HW=3};
\draw[-{Stealth},thick] (follower.north) to[bend left=18] node[above]{REPLICA\_FETCH} (leader.north);
\draw[-{Stealth},thick] (leader.south) -- node[left]{report LEO} (auth.north west);
\draw[-{Stealth},thick] (follower.south) -- node[right]{report LEO} (auth.north east);
\end{tikzpicture}
\caption{Replica logs are fetched over TCP, while progress is serialized by a trusted authority. In this example, HW=min(5,3)=3.}
\end{figure}

\lessonsection{21}{Fetch-Based Replication}
\goals{Make a second log a verifiable disk replica: a follower fetches leader records over the network and appends them to its own log.}{Understand why replicas have separate logs, why replication reads may pass the ordinary-consumer visibility boundary, and why one poll does not prove synchronization is complete.}
\background{Prerequisite route: the course sequence requires\stepref{20}{Consumer-Group Fencing}; implementation also uses the \pathcode{PartitionLog} from\stepref{6}{Segmented Logs} and request/reply handling from\stepref{12}{Broker Requests}. A leader is the broker currently accepting writes for a partition; a follower stores another copy of that partition's log and actively fetches records from the leader. Each broker has its own durable log; replication is not shared objects or copied memory. LEO is the next offset in a local log; HW is the exclusive boundary ordinary consumers may read. A follower needs the leader's uncommitted tail to catch up, so \pathcode{REPLICA_FETCH} can read through leader LEO, whereas ordinary consumer \pathcode{FETCH} reads only below HW. One poll completes only “fetch and append locally”; Step 22 reports the successfully persisted LEO to the authority.}
\providededit{The three-broker TCP harness, \pathcode{Messages.ReplicaFetchRequest} / \pathcode{ReplicaFetchBody}, \pathcode{ReplicaState}, \pathcode{PartitionLog} lifecycle, and framed RPC are provided. The initial authority reportedLEO is 0; do not mistake a direct write to the leader disk for a reported replica position.}{\pathcode{exercises/src/main/java/io/simplekafka/replication/FollowerReplicator.java}: \method{pollOnce(int maxRecords,int maxBytes)} returns the number appended in this poll; also connect the \pathcode{REPLICA_FETCH} branch in \method{handle(Messages.Request)} in \pathcode{exercises/src/main/java/io/simplekafka/broker/BrokerHandler.java}.}
\paragraph{Inputs, outputs, and state.}\pathcode{maxRecords} is a positive record-count limit and \pathcode{maxBytes} is a positive encoded-byte budget; records are indivisible, so a budget that cannot fit the next whole record may make no progress. The constructor requires a nonnegative broker id matching \pathcode{ReplicaState}; the topic-partition, log, borrowed RpcClient, and state are non-null. The poll returns the number actually appended; an empty reply returns 0. The caller owns the RPC client, which the poll does not close. Success changes follower disk and local LEO; it does not change authority reportedLEO, ISR, or HW.
\lessoncontract{First verify under the state monitor that the local broker is online and a follower, and capture its epoch; issue the request outside the monitor; validate the response; then recheck epoch/role and append under the same monitor. The request carries the current partition, local LEO, broker id, captured epoch, both positive budgets, and \pathcode{recoveryRead=false}.}
\begin{itemize}
\item The response must be a replica-fetch body with the captured epoch; its leader LEO cannot be below local LEO. Returned records must start at local LEO, have contiguous offsets, and be below the body's leader LEO. Receiving bytes is not success: only after \pathcode{PartitionLog.append} succeeds can the method count those records.
\item A nonpositive limit or unexpected body is \method{INVALID_REQUEST}; offline/already-leader local role is \method{NOT_LEADER}; a local epoch change during RPC or a wrong response epoch is \method{FENCED_EPOCH}; a regressing leader LEO, missing offset, wrong first offset, or misaligned append range is \method{CORRUPT_RECORD}. A remote leader error (for example \method{OFFSET_OUT_OF_RANGE} below its retained start) and its message are propagated; a local storage error retains its error.
\item None of these validation failures may change the follower log; if its role changes during the network request, it must not append. An empty reply at leader LEO is valid; repeated empty polls do not change file bytes. The broker caps an oversized positive byte request to one reply-frame budget; later polls copy the remaining contiguous suffix. If local writing returns \method{STORAGE_ERROR}, do not count those records or report progress; some disk side effect may already have happened, so do not assume rollback and follow the log's failed/reopen recovery path. This step does not call report or advance HW.
\end{itemize}
\lessonhints{Track “network reply received” and “follower disk append succeeded” separately.}{Compare response epoch, first offset, contiguous offsets, and leader LEO; do not substitute HW for LEO.}{Poll the same follower again, then supply a wrong epoch, wrong first offset, and an offset beyond leader LEO; verify that its LEO and file bytes remain unchanged.}
\expected{Worked example assumptions: partition start 0, epoch 0, broker 1 is leader and broker 2 follower; the leader disk has three ordinary records with null keys, one-byte ASCII values \pathcode{a,b,c}, and timestamps 11, 12, 13, so each occupies 33 bytes. Leader LEO=3, HW=0, and initial authority reportedLEO is still 0 for all three brokers. Broker 2 calls \pathcode{pollOnce(10,16384)}, receives offsets 0–2, appends all three, returns 3, and has local LEO=3; authority is unchanged. A second poll receives an empty reply, returns 0, and leaves file bytes unchanged. Ordinary consumers cannot yet see these records at HW=0, while replica requests can read them.}
\paragraph{Tests and completion.}The real \method{Step21Test} covers \method{replicaFetchIncludesUncommittedLeaderTailAndReturnsEmptyAtLeo}, \method{followerCopySurvivesEmptyPollsAndCatalogReopen}, \method{boundsLargeReplicaFetchesAndCopiesTheCompleteSuffix}, \method{rejectedReplicaFetchRequestsPreserveDiskAndAuthority}, \method{followerAheadOfLeaderIsRejectedWithoutTruncatingItsLog}, and \method{malformedRealTcpReplicaFetchRepliesAreRejectedBeforeAppend}: real TCP, disk bytes, empty polls, frame-bounded paging, no truncation when a follower is ahead, side-effect-free epoch/offset errors, and reopened contents. The single-step command diagnoses Step21; the cumulative command runs Steps1–21. An unfinished skeleton is expected to fail, and a green single-step test does not complete Chapter 6.
\pitfalls{Sharing the leader's \pathcode{PartitionLog} bypasses the network and follower persistence; treating HW as leader LEO prevents copying the uncommitted tail. Real Kafka's replica protocol is more involved; \pathcode{REPLICA_FETCH} is the course protocol, not evidence about Kafka's internal frames or control plane.}
\lessoncommand{21}
\lessonquestions{Why can a replica fetch obtain records above HW that ordinary consumers cannot yet see?}{Why must the follower not report its new LEO to the authority before its local append succeeds?}
\paragraph{Self-check answers.}Replication must transfer the leader's uncommitted tail so the follower can catch up to the current LEO; ordinary consumers use a different HW boundary. Persisting before reporting prevents the authority from counting data that is not yet on the follower's disk.

\lessonsection{22}{ISR and High Watermark}
\goals{Calculate the consumer-visible boundary only from validated progress reports, and remove long-lagging followers under a controlled clock.}{Understand that physical LEO, authority reports, ISR, minISR, and HW are different state.}
\background{Prerequisite:\stepref{21}{Fetch-Based Replication}. A follower may have written data in its own log without reporting it yet. The authority keeps each broker's most recently successful reported LEO for a partition; this can temporarily differ from the physical LEO just read from that log. ISR is the set of replicas currently participating in acknowledgment; minISR is the minimum ISR size for ALL requests; HW is the exclusive end of the prefix they have confirmed together. Only a successful report changes authority progress.}
\providededit{Per-partition locks and Conditions, assigned-replica membership, online status, authority progress containers, and an injectable millisecond \pathcode{TimeSource} are provided.}{\pathcode{exercises/src/main/java/io/simplekafka/replication/ReplicationTracker.java}: \method{report(int brokerId,int epoch,long leo)} returns void; \method{expireLagging()} returns the set of removed broker ids; \method{advanceHighWatermark()} returns the current HW.}
\paragraph{Inputs, outputs, and state.}\pathcode{report} requires a nonnegative broker id belonging to the assignment and online, plus nonnegative epoch/LEO and the current leader epoch; invalid input leaves the snapshot unchanged. A leader report cannot move backward; a follower report cannot exceed the leader's reported LEO or move backward in the same epoch. A successful report updates that broker's reportedLEO; a follower refreshes its last-caught-up time only when its LEO reaches the leader LEO observed for that report. A report does not readmit a member already removed from ISR.
\lessoncontract{If ISR contains the leader, all member progress exists, and ISR size is at least minISR, the candidate HW is the minimum latest successful reported LEO in ISR, applied as \pathcode{HW=max(oldHW,candidate)}; otherwise HW freezes. HW never decreases.}
\begin{itemize}
\item Negative broker id, epoch, or LEO; unknown/unassigned/offline broker; a backward leader LEO; or a follower ahead of or behind its permitted progress is \method{INVALID_REQUEST}. A wrong epoch is \method{FENCED_EPOCH}. A rejected report changes neither reportedLEO, ISR, nor HW.
\item \pathcode{expireLagging()} checks non-leaders in ISR. A member with no catch-up time or with \pathcode{now-lastCaughtUp >= lagTimeoutMillis} is removed and included in the immutable returned set. Timeout 0 means immediate expiration; if now precedes the catch-up time, elapsed time does not expire it. Expiration removes neither assignment nor log, and never removes the leader.
\item After removals, reconsider HW; if the remaining ISR is below minISR, it freezes. Successful reports, actual expiration, or HW advancement notify waiters; \pathcode{advanceHighWatermark()} recomputes HW from the current ISR and returns it. If HW actually advances, it signals waiters; the method itself does not report replica progress or expire lagging replicas.
\end{itemize}
\lessonhints{Validate follower reports against “not beyond leader” and “not backward” before calculating HW.}{If minISR is insufficient, do not advance from the remaining LEOs; refresh catch-up time only on catch-up.}{Use a fixed clock to check one millisecond before a deadline, the exact deadline, and one invalid report; compare snapshots without sleeps.}
\expected{Worked progress example: replicas/ISR={1,2,3}, minISR=2, previous HW=0, and reported LEOs for leader1 and followers2,3 are 5,3,4. The minimum candidate is3, so HW=3. Follower2 reports5: candidate and HW become4. Follower3 reports5: HW becomes5. Separately, follower3 last caught up at clock=1000ms with lag timeout=1000ms; now=1999ms retains it, while now=2000ms expires it at the exact boundary. Its log is not rewritten.}
\paragraph{Tests and completion.}\method{Step22Test.rejectsInvalidProgressAndExpiresOnlyReplicasThatMissTheExactCatchUpDeadline} starts the clock at 1000 and verifies that negative progress, a future epoch, unassigned/offline replicas, and follower overrun/rollback leave the snapshot unchanged; no member expires at 1999, while broker 3 expires at 2000. It also checks HW freeze as ISR shrinks, progress to 4 after a new report, and that HW stays at 4 when only the leader reports 5. The single-step command diagnoses Step 22; the cumulative command runs Steps 1–22. A green single-step test does not complete Chapter 6.
\pitfalls{Online does not mean ISR, and physical log LEO does not mean authority-known LEO; using the former directly for HW could expose unconfirmed data. The course serializes state in a trusted single-JVM authority; real Kafka does not guarantee the global HW monotonicity across epochs used here.}
\lessoncommand{22}
\lessonquestions{Why can insufficient minISR prevent HW from advancing even when some replicas have written the data?}{Why does an online follower that has not caught up to leader LEO fail to refresh its catch-up time?}
\paragraph{Self-check answers.}HW marks a readable common prefix, not one replica's write position; an undersized ISR cannot meet this course's acknowledgment condition. Only a report that has caught up proves the follower reached the current leader tail; a live connection alone proves nothing about its data.

\lessonsection{23}{Produce Acknowledgments}
\goals{Explain the difference among a write, a replication commit, and an acknowledgment delivered to the caller, then implement both wait rules.}{Understand half-open offset ranges, Conditions, monotonic-clock deadlines, and why a timeout does not undo an append.}
\background{Prerequisite:\stepref{22}{ISR and High Watermark}. An ACK setting tells the producer how far the broker must wait before replying: \pathcode{LEADER} returns when the leader's local append completes; \pathcode{ALL} requires sufficient minISR before append and waits for HW to reach the batch's \pathcode{nextOffset}. For a returned range [a,b), HW\(\ge b\) means the entire batch is in the committed prefix; readers can fetch only offsets below HW. A deadline is an absolute instant in the \pathcode{System.nanoTime()} time domain, while timeoutMillis is a request duration; they are not interchangeable.}
\providededit{The ReplicationTracker partition lock/Condition, \pathcode{AppendResult}, \pathcode{Acks}, \pathcode{RecordData}, and \pathcode{TimeSource} are provided; \pathcode{deadlineAfterMillis(long)} is provided too.}{\pathcode{exercises/src/main/java/io/simplekafka/replication/AckPolicy.java}: \method{validateBeforeAppend(Acks acks)} returns void; \method{await(AppendResult result,Acks acks,int epoch,long deadlineNanos)} returns void.}
\paragraph{Inputs, outputs, and side effects.}\pathcode{validateBeforeAppend} accepts non-null LEADER/ALL; LEADER needs no ISR check, while ALL checks ISR size before append. \pathcode{await} accepts a non-null result and ACK mode, the epoch captured for that append, and an absolute monotonic deadline; LEADER returns immediately, while ALL waits for the result's \pathcode{nextOffset}. The caller's \pathcode{ReplicatedPartition.produce} rejects negative timeout or epoch. Null arguments throw NullPointerException under the interface contract. Any wait error after append does not retract the disk record.
\lessoncontract{Before append, ALL with ISR<minISR reports \method{NOT_ENOUGH_REPLICAS}, so no record has yet been appended. After append it rechecks in this order: \pathcode{epoch → ISR → HW → deadline}. An epoch change reports \method{FENCED_EPOCH}; ISR below minISR reports \method{NOT_ENOUGH_REPLICAS}; if HW has reached result nextOffset, return success; otherwise an expired deadline reports \method{REQUEST_TIMEOUT}. Satisfied HW takes precedence over an expired deadline. Condition waiting releases the partition lock; after waking, recheck all conditions. Interruption restores the thread's interrupt flag and reports \method{REQUEST_TIMEOUT}, retaining the interrupt as cause. LEADER's await completes immediately.}
\lessonhints{Separate the pre-append minISR check from the post-append HW wait; their failures have different side effects.}{A Condition signal only says state may have changed; reacquire the lock and check every condition in order.}{Use latches to control report order and verify that a report wakes ALL; then test zero timeout, epoch change, ISR reduction, and interruption against LEO/HW.}
\expected{Assume leader epoch=0, ISR={1,2,3}, minISR=2. Append one ordinary record with null key and empty value (32 encoded bytes), returning [0,1); leader LEO=1, follower LEO=0, HW=0. ALL with timeout 0 reports \method{REQUEST_TIMEOUT} after the append: LEO stays 1, HW stays 0, the caller has no receipt, and a consumer cannot see offset 0. Later both followers copy and report LEO=1, so HW may reach 1 and the record becomes visible; the old request does not retroactively receive a successful receipt. If ISR was already too small before append, ALL instead reports \method{NOT_ENOUGH_REPLICAS} with LEO still 0.}
\paragraph{Tests and completion.}\method{Step23Test} contains \method{insufficientIsrRejectsAllBeforeAppendWhileLeaderAckStillAppends}, \method{exactHighWatermarkBoundarySucceedsButOneOffsetBelowWaitsForProgress}, \method{zeroTimeoutRetainsTheAppendWithoutAdvancingTheHighWatermark}, \method{epochChangeWhileWaitingFencesTheAcknowledgement}, and \method{isrReductionWhileWaitingFailsAndInterruptionPreservesTheInterruptFlag}. They verify side-effect-free prechecks, HW boundary, retained timeout append, epoch/ISR wake-ups, and interrupt preservation. The single-step command diagnoses Step 23; the cumulative command runs Steps 1–23. Chapter completion requires cumulative Step 24.
\pitfalls{Do not wait on replication while holding the partition lock or the only progress-reporting thread may be blocked; wall-clock adjustments do not belong in monotonic deadlines. A timeout means the caller did not receive confirmation by the deadline, not that no append occurred. The course supports only two ACK modes and does not retry for the caller.}
\lessoncommand{23}
\lessonquestions{Why must epoch, ISR, HW, and deadline be checked again after a Condition releases the lock?}{After an ALL timeout, what can be true about LEO, HW, and the caller's knowledge of the append?}
\paragraph{Self-check answers.}Any state may change while waiting, and a signal does not prove the condition is satisfied. LEO may have advanced while HW has not, so data can become visible later even though the caller never received a safe success receipt.

\lessonsection{24}{Concurrent Replicated Broker}
\goals{Combine leader validation, appends, replication acknowledgments, and ordinary reads in a broker serving concurrent requests.}{Understand authority/local-log lock order, HW read clipping, and socket/Condition lifecycle.}
\background{Prerequisites are\stepref{21}{Fetch-Based Replication},\stepref{22}{ISR and High Watermark}, and\stepref{23}{Produce Acknowledgments}. An ordinary consumer can read only the complete committed prefix offset<HW; a valid offset within the retained range but \pathcode{offset >= HW} returns an empty list. An offset beyond LEO is out of range. Replica fetch is a separate request that may obtain the uncommitted tail. Produce permission checks, append, and leader progress are linearized under the authority per-partition lock; the lock order is authority lock→local \pathcode{PartitionLog} monitor. No network call may hold either lock; an ALL wait must run outside them so follower reports can proceed.}
\providededit{Independent TCP endpoints, broker workers, ReplicaState, ReplicationTracker, AckPolicy, PartitionLog, and broker routing are provided.}{\pathcode{exercises/src/main/java/io/simplekafka/replication/ReplicatedPartition.java}: \method{produce(List<RecordData> records,Acks acks,int epoch,long timeoutMillis)} returns \method{AppendResult}; \method{fetch(long offset,int maxRecords,int maxBytes,int epoch)} returns \pathcode{List<LogRecord>}.}
\paragraph{Parameters, results, and state.}\pathcode{produce} requires non-null records/ACK mode, nonnegative epoch and timeoutMillis, and a currently online leader matching the request epoch; success returns a half-open append range and may append to the log/update replication progress. \pathcode{fetch} requires nonnegative offset/epoch and positive record/byte limits; it only reads log and authority state, returning complete records bounded by count, encoded bytes, and HW. It never splits a record. The provided handler routes network workers.
\lessoncontract{After argument/null checks, produce first runs ALL's pre-append minISR check, before online/role/epoch checks. Under the authority lock it then synchronizes and validates online/leader/epoch, rechecks ALL minISR to cover state changes since the first check, appends locally, reports leader progress/possible HW advancement, unlocks, and waits for the ACK. Either failed minISR check or role/epoch failure must not append; the first check can make a minISR error take precedence over \method{NOT_LEADER}/\method{FENCED_EPOCH}. Negative parameters or an empty produce batch are \method{INVALID_REQUEST}, a non-leader is \method{NOT_LEADER}, a stale epoch is \method{FENCED_EPOCH}, and local storage failure is \method{STORAGE_ERROR}. Post-append \method{NOT_ENOUGH_REPLICAS}, \method{REQUEST_TIMEOUT}, or \method{FENCED_EPOCH} does not roll back the append.}
\begin{itemize}
\item Fetch rejects a negative offset/epoch or nonpositive budget with \method{INVALID_REQUEST}; a stale epoch is \method{FENCED_EPOCH}; a requester that is not the online leader is \method{NOT_LEADER}; offset<logStart or offset>LEO is \method{OFFSET_OUT_OF_RANGE}.
\item Fetch's legal offset domain is the closed interval [logStart,LEO]; a value in that domain with offset\(\ge HW\) returns empty, while offset<HW returns only whole records with offset<HW. LEO is the next offset, so offset=LEO can be valid and empty; greater than LEO is invalid.
\item Keep lock order authority→local log; do not reverse it or hold locks across a broker RPC. Broker close must let sockets and waiting ACK workers finish; a pending append may remain in the log.
\end{itemize}
\lessonhints{Draw the lock intervals for produce, replica poll, and consumer fetch, placing every TCP call outside them.}{Check that HW clipping happens before an ordinary fetch reply, rather than asking consumers to discard uncommitted records.}{Pause replication after an ALL append, fetch from another connection, and issue a LEADER append concurrently; resume replication and verify all three can finish.}
\expected{Concurrency example assumptions: leader=1, followers=2/3, epoch=0, initially empty log. First a LEADER append writes offset 0 and is replicated; all three LEOs=1 and HW=1. Then ALL appends offset 1 and waits; leader LEO=2, HW=1. A separate fetch(0) sees only offset 0, and fetch(1) is empty. During the wait, LEADER can append offset 2 (range [2,3)), but fetch still sees only offset 0. The followers then copy offsets 1–2 and report LEO=3; HW advances to 3, ALL returns [1,2), and ordinary fetch can now see offsets 0, 1, and 2. For the byte-budget fixture, assume each ordinary record has null key and empty value, hence 32 encoded bytes: among four appended records with HW=2, fetch from 0 with maxBytes=31 returns none, 32 or 63 returns offset 0 only, and 64 returns offsets 0 and 1; maxRecords=1 still caps the result at one record.}
\paragraph{Tests and completion.}\method{Step24Test.pendingAllAckLeavesFetchAtOldHwAllowsLeaderProduceAndCompletesAfterTcpReplication} checks reads at old HW and progress of another produce while waiting; \method{fetchHonorsHwLeoAndWholeRecordByteBudgetsAndRejectsInvalidRequests} checks the 31/32/63/64-byte boundaries, HW/LEO, and out-of-range case; \method{allPrecheckDoesNotAppendAndZeroTimeoutRetainsAnInvisibleAppendUntilReplication} and \method{stoppingLeaderCompletesAnOutstandingAllAckRequestWithoutLeakingTheProducer} check failure side effects and shutdown wake-up. The single-step command diagnoses Step 24; the cumulative command runs Steps 1–24 and is the Chapter 6 completion boundary.
\pitfalls{Holding a lock across a socket can prevent the only replication-progress thread from running; HW is not a group committed offset. The three independent logs replicate over real TCP, but the authority remains a single point in one shared JVM, so this is not production HA.}
\lessoncommand{24}
\lessonquestions{Why can another fetch safely read the old HW prefix while ALL is waiting?}{Why might closing only the broker socket be insufficient to end every waiting thread?}
\paragraph{Self-check answers.}Fetch reads the already committed, immutable prefix and cannot leak the uncommitted offset. A Condition waiter also needs a wake-up and an exit check on state change/shutdown; closing a connection alone does not replace worker and waiter lifecycle management.
